\section{The Ising ferromagnet in two dimensions}

Throughout, $S = \mathbb{Z}^2$, $E = \{\pm 1\}$ (Lean: \texttt{Bool}), and $\gamma^{\beta}$ is the
Gibbsian specification of the nearest-neighbour Ising potential with coupling $J = 1$ and external
field $h = 0$ at inverse temperature $\beta$.

\subsection*{The Peierls argument}

\begin{lemma}[The Peierls estimate]
    \label{lem:ising-peierls}
    \uses{def:spec}
    \lean{MeasureTheory.GibbsMeasure.Peierls.r,
      MeasureTheory.GibbsMeasure.Peierls.isingSpecification_cube_eq_false_le,
      MeasureTheory.GibbsMeasure.Peierls.r_le_quarter,
      MeasureTheory.GibbsMeasure.tendsto_r_atTop}
    \leanok

    In a cube $\Lambda = [-N,N]^2$ with $+$ boundary condition,
    \[\gamma^{\beta}_{\Lambda}(\sigma_a = -1 \mid +) \;\leq\; r(\beta)\]
    for every site $a$, where $r(\beta) = \sum_{\ell \geq 1} \ell\, 4096^{\ell} e^{-2\beta\ell}$ is
    the Peierls series: a minus spin at $a$ forces a contour around $a$, and each contour of length
    $\ell$ costs $e^{-2\beta\ell}$. Moreover $r(\beta) \leq \tfrac14$ once
    $\beta \geq 8\log 2$, and $r(\beta) \to 0$ as $\beta \to \infty$.

    The chain from the cube estimate to the phase transition is proved once, for an arbitrary
    bound $\rho$ dominating the cube estimate
    (\texttt{exists\_two\_shiftInvariant\_gibbs\_of\_cube}), and instantiated at $r$ and at the
    sharpened bound $r'$ of the refined contour count.
\end{lemma}

\begin{theorem}[Georgii Theorem (6.9), spontaneous magnetisation]
    \label{thm:ising-69-two-phases}
    \uses{lem:ising-peierls, def:gibbs-meas}
    \lean{MeasureTheory.GibbsMeasure.Peierls.exists_plus_phase,
      MeasureTheory.GibbsMeasure.Peierls.exists_two_shiftInvariant_gibbs,
      MeasureTheory.GibbsMeasure.Peierls.exists_two_shiftInvariant_gibbs_of_large_beta,
      MeasureTheory.GibbsMeasure.Peierls.nontrivial_GP_isingSpecification_of_large_beta,
      MeasureTheory.GibbsMeasure.PeierlsSharp.exists_two_shiftInvariant_gibbs_sharp}
    \leanok

    For $\beta \geq 8 \log 2$ — and, with the sharpened contour count, already for
    $\beta \geq \tfrac12\log 9 = \log 3$ — the two-dimensional Ising ferromagnet has two distinct
    shift-invariant Gibbs measures $\mu_\pm$, exchanged by the spin flip, with
    \[\mu_+(\sigma_0 = -1) \;<\; \tfrac12 \;<\; \mu_-(\sigma_0 = -1).\]
    The plus phase arises as a cluster point of the cube averages with $+$ boundary condition;
    shift-invariance comes from averaging over cube translates.
\end{theorem}

\begin{theorem}[Georgii Theorem (6.9), the low-temperature limit]
    \label{thm:ising-69-limit}
    \uses{thm:ising-69-two-phases}
    \lean{MeasureTheory.GibbsMeasure.Peierls.shiftInvariantGibbs,
      MeasureTheory.localDist, MeasureTheory.localDistSet,
      MeasureTheory.GibbsMeasure.Peierls.tendsto_localDistSet_shiftInvariantGibbs_dirac}
    \leanok

    With $d$ a metric for the topology of local convergence (Georgii Remark (4.3)(3), the explicit
    metric of \texttt{Topology/LocalMetric.lean}),
    \[\lim_{\beta \to \infty} d\bigl(\mathcal{G}_{\Theta}(\beta\Phi), \delta_+\bigr)
      = \lim_{\beta \to \infty} d\bigl(\mathcal{G}_{\Theta}(\beta\Phi), \delta_-\bigr) = 0:\]
    at low temperature the shift-invariant Gibbs measures approach the two ground states.
\end{theorem}

\subsection*{Griffiths' inequality and the spontaneous magnetisation}

\begin{lemma}[GKS and the finite-volume magnetisation]
    \label{lem:ising-gks-fvmag}
    \lean{MeasureTheory.GibbsMeasure.fvMag,
      MeasureTheory.GibbsMeasure.fvMag_eq_corr,
      MeasureTheory.GibbsMeasure.fvMag_nonneg,
      MeasureTheory.GibbsMeasure.fvMag_mono,
      MeasureTheory.GibbsMeasure.plusMag,
      MeasureTheory.GibbsMeasure.tendsto_fvMag,
      MeasureTheory.GibbsMeasure.spontaneousMagnetisation}
    \leanok

    The finite-volume $+$-boundary magnetisation
    $\langle\sigma_i\rangle_{\Lambda}^{+}(\beta) = 2\gamma^{\beta}_{\Lambda}(\sigma_i = 1 \mid +)-1$
    is a GKS correlation of a ferromagnetic interaction; by Griffiths' inequalities it is
    nonnegative and nondecreasing in $\beta$, and it converges along the volumes to
    $\mu_+^{\beta}(\sigma_i)$. The \textbf{spontaneous magnetisation} is
    $m^*(\beta) = \mu_+^{\beta}(\sigma_0)$.
\end{lemma}

\begin{theorem}[Lebowitz--Martin-L\"of/Ruelle]
    \label{thm:ising-lml}
    \uses{lem:ising-gks-fvmag}
    \lean{MeasureTheory.GibbsMeasure.nontrivial_GP_ising2D_iff_spontaneousMagnetisation_pos,
      MeasureTheory.GibbsMeasure.nontrivial_GP_ising2D_of_nontrivial_of_le,
      MeasureTheory.GibbsMeasure.isingNonUniqueness_eq_setOf_spontaneousMagnetisation_pos}
    \leanok

    For $\beta \geq 0$, the two-dimensional Ising ferromagnet has more than one Gibbs measure if
    and only if $m^*(\beta) > 0$. Consequently non-uniqueness is an upper set in $\beta$:
    Georgii's route, via Griffiths, to the existence of a critical inverse temperature.
\end{theorem}

\subsection*{The critical inverse temperature}

\begin{definition}[The critical inverse temperature]
    \label{def:ising-betac}
    \uses{thm:ising-lml}
    \lean{MeasureTheory.GibbsMeasure.isingNonUniqueness,
      MeasureTheory.GibbsMeasure.isingBetaC,
      MeasureTheory.GibbsMeasure.betaCOf,
      MeasureTheory.GibbsMeasure.isingBetaC_eq_betaCOf_spontaneousMagnetisation}
    \leanok

    $\beta_c = \inf\set{\beta \geq 0 : |\mathcal{G}(\beta\Phi)| > 1}$. By the
    Lebowitz--Martin-L\"of/Ruelle criterion this is the same number as the threshold
    $\inf\set{\beta \geq 0 : m^*(\beta) > 0}$ of the order parameter, which is how Georgii defines
    it.
\end{definition}

\begin{theorem}[The sharp dichotomy at $\beta_c$]
    \label{thm:ising-betac-dichotomy}
    \uses{def:ising-betac}
    \lean{MeasureTheory.GibbsMeasure.existsUnique_of_lt_isingBetaC,
      MeasureTheory.GibbsMeasure.nontrivial_of_isingBetaC_lt,
      MeasureTheory.GibbsMeasure.isUpperSet_isingNonUniqueness,
      MeasureTheory.GibbsMeasure.ising_critical_temperature,
      MeasureTheory.GibbsMeasure.ising_sharp_phase_transition,
      MeasureTheory.GibbsMeasure.spontaneousMagnetisation_eq_zero_of_lt_isingBetaC,
      MeasureTheory.GibbsMeasure.spontaneousMagnetisation_pos_of_isingBetaC_lt}
    \leanok

    For $0 \leq \beta < \beta_c$ the Gibbs measure is unique and $m^*(\beta) = 0$; for
    $\beta > \beta_c$ there are several Gibbs measures and $m^*(\beta) > 0$ — spontaneous
    magnetisation. This is Georgii's statement in §6.2 after (6.9), for the genuine order
    parameter, not an abstract stand-in.
\end{theorem}

\begin{theorem}[Bounds on $\beta_c$]
    \label{thm:ising-betac-bounds}
    \uses{def:ising-betac}
    \lean{MeasureTheory.GibbsMeasure.artanh_le_isingBetaC,
      MeasureTheory.GibbsMeasure.le_isingBetaC,
      MeasureTheory.GibbsMeasure.isingBetaC_le,
      MeasureTheory.GibbsMeasure.isingBetaC_mem_Icc_artanh,
      MeasureTheory.GibbsMeasure.isingBetaC_pos,
      MeasureTheory.GibbsMeasure.Dobrushin.isDobrushin_isingSpecification,
      MeasureTheory.GibbsMeasure.Dobrushin.isDobrushin_isingSpecification_tanh}
    \leanok

    \[\operatorname{artanh}\tfrac14 \;\leq\; \beta_c \;\leq\; \log 3 .\]
    The lower bound is Dobrushin uniqueness with the sharpened single-site bound
    $2d\tanh|\beta J| < 1$ of (8.9)(2)/(8.10); the upper bound is the sharpened Peierls
    threshold $\log 3$. (Onsager's exact
    $\beta_c = \tfrac12\log(1 + \sqrt2) \approx 0.4407$ is not proved in Georgii and not claimed
    here; $\operatorname{artanh}\tfrac14 \approx 0.2554$.)
\end{theorem}
